\documentclass[letterpaper,10pt,conference]{ieeeconf}

\IEEEoverridecommandlockouts
\overrideIEEEmargins

% math + refs
\usepackage{amsmath,amssymb,amsfonts,mathtools}
\usepackage[table]{xcolor}
\newcommand{\chg}[1]{\textcolor{red}{#1}}

% (optional but common)
\usepackage{graphicx}
\usepackage{subcaption}
\usepackage{cite}
\usepackage{hyperref}
\usepackage{amsmath,amssymb}
\usepackage{algorithm}
\usepackage{algpseudocode}
\usepackage{booktabs}

\begin{document}

\title{\LARGE \bf
Scenario-based Safe Horizon MPCC with Wasserstein-DRO Mode Reweighting and Scenario Injection for Targets with Switching Dynamics
}

\author{Stephen Crawford%
\thanks{Department of Computer Science, Brown University, Providence, RI, USA.
{\tt\small stephen\_crawford@brown.edu}}%
\and
Nora Ayanian%
\thanks{Department of Computer Science, Brown University, Providence, RI, USA.
{\tt\small nora\_ayanian@brown.edu}}%
}

\maketitle

\begin{abstract}
Scenario-based MPC offers finite-sample guarantees for chance-constrained planning under i.i.d.\ uncertainty samples, but switching behaviors can cause rare yet safety-critical modes to be under-represented at realistic scenario budgets. We present a Safe Horizon Model Predictive Contouring Control (SH-MPCC) framework coupled with a Wasserstein distributionally robust optimization (WDRO) mechanism. Given a nominal mode belief $\widehat p_t^v$, WDRO produces a high-danger distribution $q^\star$ in a Wasserstein ball that redistributes weight to modes with larger predicted safe-horizon clearance violations. \chg{The Wasserstein radius is calibrated as a finite-sample $W_1$ concentration bound, so the ambiguity set covers the true mode law with high confidence and contracts to the nominal belief as observations accumulate.}
We study two uses of $q^\star$: a sampling strategy, which samples modes from $q^\star$ to improve finite-sample mode coverage under switching, and scenario injection, which injects $K$ deterministic high-risk mode constraints as an additional constraint layer. These strategies are tested across a number of environements and scenario budgets and compared against base SH-MPCC. Across two environments, WDRO-sampling reduces missed-mode rates by approximately 10 percentage points relative to risk-softmax sampling at comparable collision rates. For injection, WDRO provides its largest safety gains when multiple distinct risky behaviors must be represented under a finite injection budget ($K>1$), improving mode-set selection. A history-length sweep further shows WDRO-injection is the most stable method under limited mode observations, and ground-cost ablations indicate performance is not sensitive to the exact transport cost choice. Overall, WDRO mode reweighting and lightweight injection provide a practical mechanism for improved coverage and safety in SH-MPCC with low computational overhead.
\end{abstract}

%==============================================================================
% SECTION I: INTRODUCTION
%==============================================================================

\section{Introduction}

Autonomous systems operating in shared environments must make decisions under uncertainty caused by other dynamic agents.
Even when the ego dynamics are well modeled, safe motion planning becomes difficult when nearby agents exhibit switching behaviors.
This is commonly seen in the context of driving where human drivers infer intent from subtle cues and remain cautious about rare but plausible actions (e.g., an unexpected turn or lane change). This has motivated extensive work on behavior modeling and perception \cite{Ha2013pedestrianBehavior,He1995socialForce,Zi2015youngOldCrossing,Hu2017bikeSocial,Hu2017bikeBehavior}.

A broad set of approaches for planning under uncertainty falls under robust and stochastic model predictive control (RMPC/SMPC).
RMPC addresses bounded (set-membership) uncertainty, while SMPC incorporates probabilistic uncertainty via chance constraints \cite{Mesbah2016SMPCOverview}.
Chance-constrained MPC offers a principled way to minimize risk at the cost of performance, but exact chance constraints are typically intractable in nonlinear, multi-agent settings \cite{Ch1959CCP,schwarm1999chance}.

A key limitation of many existing methods is that they control collision risk either per time step or per obstacle.
Per-time-step approaches often use Boole's inequality across the horizon, yielding conservative risk allocation.
Per-obstacle approaches do not directly capture the joint event of colliding with any obstacle at any step.
In contrast, Safe Horizon MPC (SH-MPC) \cite{degroot2023SBMotionBounded} constrains the joint collision probability over the planning horizon.
Scenario optimization \cite{Cal2006ScenarioApproach,Ca2018ScenarioTheory} provides a surrogate by sampling finitely many obstacle-trajectory realizations and enforcing collision avoidance on each sample. This results in a tractable scenario program with probabilistic feasibility guarantees.

Despite these benefits, scenario-optimization based MPC approaches remain only as reliable as its sampling distribution.
When obstacles switch between modes (e.g., ``go straight'' versus ``turn''), sampling from the currently dominant mode can under-represent rare transitions.
At practical scenario budgets, safety-critical behaviors may be missed, leading to optimistic collision-risk estimates.
This issue worsens as trajectory dimension increases and as more complex agents are considered.

To address under-sampling of rare switching behaviors at practical scenario budgets, we focus on improving which modes are represented in the scenario set rather than increasing the number of sampled trajectories.
We introduce a Wasserstein-DRO mode reweighting strategy that computes a high-danger distribution $q^\star$ in a Wasserstein ball around a nominal belief $\widehat p_t^v$. 
We study two complementary mechanisms: sampling scenarios from $q^\star$, and injecting an additional deterministic scenario corresponding to the highest-risk mode under $q^\star$. 
We pair these mechanisms with SH-MPCC (SH-MPC with contouring control), which provides the primary collision avoidance mechanism.

Our paper contributes: (i) a WDRO mode reweighting rule for SH-MPCC under switching, (ii) empirical evidence that WDRO-sampling improves finite-sample mode coverage relative to risk-softmax sampling, (iii) a mode-set injection strategy where WDRO is demonstated to inject better modes than a greedy risk baseline, and (iv) robustness results across observation regimes and ground-cost choices with low runtime overhead.

The remainder of this paper is organized as follows.
Section~\ref{related} reviews related work.
Section~\ref{preliminaries} introduces scenario MPC preliminaries and our Wasserstein-DRO reweighting and injection mechanisms.
Section~\ref{methods} presents the methodology.
Section~\ref{results} reports results and ablations.
Section~\ref{conclusion} concludes.

%==============================================================================
% SECTION II: RELATED WORK
%==============================================================================

\section{Related Work}\label{related}

Scenario-based optimization provides a tractable surrogate for chance-constrained control by enforcing constraints on finitely many sampled ``scenarios'' of the uncertainty \cite{Cal2006ScenarioApproach,Ca2018ScenarioTheory}. In motion planning and MPC, this approach has been used to reduce conservatism relative to convex relaxations while preserving real-time feasibility \cite{de2021ScenarioBasedTraj,degroot2023SBMotionBounded}.

De Groot et al.\ \cite{de2021ScenarioBasedTraj} control collision risk via per-step chance constraints within a receding-horizon replanning loop. The same authors \cite{degroot2023SBMotionBounded} extend this framework to bound the joint probability of collision over the trajectory, reducing conservatism while retaining scenario-based tractability; this is the Safe Horizon MPC (SH-MPC) method upon which our implementation builds. Other scenario optimization-based MPC variants introduce candidate-path representations \cite{Shang2021SceMPC} or reinforcement-learning-based components \cite{Jekl2025SceVelPlanning}, which are complementary to our focus.

More broadly, chance-constrained MPC for multi-agent collision avoidance has been studied using particle approximations, surrogate constraints, and distributional assumptions (often Gaussian) \cite{Blackmore2010ProbParticle,Gopalakrishnan2016CCMulti,Lyons2012CCMPCMulti,Zhu2019CCColl,wang2020NonGauss}. These methods typically address uncertainty under fixed probabilistic models rather than unknown switching. Approaches that explicitly model context-dependent or switching behaviors often rely on intent inference or structured prediction modules \cite{Kooij2019Context-based}.

Model Predictive Contouring Control (MPCC) has become a popular formulation for path-following, where a contouring error and lag error are penalized with respect to a reference path parametrized by a progress variable \cite{Liniger2014MPCCRacing}. Our framework embeds contouring control within the scenario-based optimization, allowing path tracking and collision avoidance to be jointly optimized.

In contrast to Wasserstein-DRO reformulations of MPC that aim to robustify the disturbance model inside the optimization \cite{kuhn2019wdro,gao2023drso,Hakobyan2020WDRO,Zhong2021DDRMPC},
we use Wasserstein-DRO as a mode reweighting and scenario selection layer.
Given a nominal belief $\widehat p_t^v$ over discrete motion modes, we compute a worst-case
distribution $q^\star$ in a Wasserstein ball and either (i) sample scenarios from $q^\star$
or (ii) inject deterministic mode-representative constraints as an additional set of constraints.
Thus, DRO is used to improve finite-sample coverage of rare but safety-critical modes at
fixed scenario budget, rather than to replace the MPC constraints with a DRO counterpart.

%==============================================================================
% SECTION III: PRELIMINARIES
%==============================================================================

\section{Preliminaries}\label{preliminaries}

\subsection{Chance-constrained MPC}\label{sec:cc_prelim}

We consider discrete-time stochastic control with state $x_t\in\mathbb{R}^n$, input $u_t\in\mathbb{U}\subset\mathbb{R}^m$, and disturbance $\omega_t$ drawn from some probability distribution $\mathbb{P}_\omega$:
\begin{equation}\label{eq:dyn_general}
x_{t+1} = f(x_t,u_t,\omega_t), \qquad \omega_t \sim \mathbb{P}_\omega,
\end{equation}
with outputs $y_t = h(x_t,u_t,\psi_t)$ and $\psi_t \sim \mathbb{P}_\psi$ \cite{Mesbah2016SMPCOverview}.
At time $t$, MPC plans over a horizon of length $N$ using predicted states $x_{t+k|t}$ for $k=0,\dots,N$ and a control sequence $U_t = (u_{t|t},\dots,u_{t+N-1|t})$.

A generic joint chance-constrained MPC problem can be written as
\begin{subequations}\label{eq:cc_mpc}
\begin{align}
\min_{U_t\in\mathbb{U}^N}\quad &
\mathbb{E}\!\left[\sum_{k=0}^{N-1} \ell(x_{t+k|t},u_{t+k|t}) + \ell_f(x_{t+N|t})\right] \\
\text{s.t.}\quad &
x_{t+k+1|t}=f(x_{t+k|t},u_{t+k|t},\omega_{t+k}),\quad x_{t|t}=x_t,\\
& \mathbb{P}\!\left(g_j(x_{t+1:t+N|t},\boldsymbol{\delta})\le 0,\ \forall j=1,\dots,s\right) \ge 1-\epsilon.
\end{align}
\end{subequations}
where $\ell$ denotes the cost function, $\ell_f$ is a terminal cost, $g_j(\cdot)$ encodes constraint violations (e.g., collision avoidance), $\boldsymbol{\delta}$ denotes the random quantities entering the constraints (e.g., obstacle trajectories), and $\epsilon$ is a defined threshold.

\subsection{Scenario approximation}

Scenario optimization replaces the chance constraint by finitely many sampled constraints \cite{Cal2006ScenarioApproach,Ca2018ScenarioTheory}. Let $\boldsymbol{\delta}\in\Delta$ denote an uncertainty instance and let $\{\boldsymbol{\delta}^{(i)}\}_{i=1}^S$ be i.i.d.\ samples. The scenario program enforces
\begin{equation}\label{eq:scenario_generic}
g(u,\boldsymbol{\delta}^{(i)}) \le 0,\quad i=1,\dots,S,
\end{equation}
where $g$ is the constraint violation function and $u$ is the decision variable. Under convexity and i.i.d.\ sampling, scenario theory provides bounds on $S$ sufficient to guarantee that the optimizer violates the original chance constraint with probability at most $\epsilon$ \cite{Cal2006ScenarioApproach,Ca2018ScenarioTheory}.

\subsection{Mode coverage metrics (diagnostic)}\label{sec:metrics_prelim}

We use two diagnostics throughout:
(i) overall missed-mode rate which is the fraction of rollouts where at least one obstacle mode that was demonstrated was not represented in the sampled scenario set within the safe horizon; and
(ii) rare-mode miss rate, which is the fraction a designated rare mode is not captured.

\subsection{Wasserstein-DRO reweighting} \label{sec:wass_prelim}

Let $\widehat p_t^v\in\Delta_M$ denote a nominal belief over the $M$ modes for obstacle $v$ at time $t$.
To reshape mode weights in a geometry-aware way, we compute a worst-case distribution $q\in\Delta_M$ in a Wasserstein ball around $\widehat p_t^v$:
\begin{equation}\label{eq:w_dro_modes}
\max_{q\in\Delta_M}\ \langle q,r\rangle
\quad\text{s.t.}\quad
W_D(q,\widehat p_t^v)\le \rho,
\end{equation}
where $r\in\mathbb{R}^M$ is a per-mode danger score computed from predicted safe-horizon clearance violations (higher indicates greater risk), $W_D$ is the discrete Wasserstein distance induced by ground cost $D$, and $\rho$ is the Wasserstein radius \cite{Esfahani2017DDRDRO,Zhong2021DDRMPC,kuhn2024WDRO}. We scale the Wasserstein ambiguity radius to reflect the reliability of the empirical mode belief.

\chg{The construction stacks two Wasserstein distances at different levels, which we distinguish explicitly to avoid conflation. The ambiguity ball is a \emph{Wasserstein-1} (Kantorovich) ball over the mode simplex $\Delta_M$: the transport cost $W_D(q,\widehat p_t^v)=\min_{\Pi\ge 0}\{\sum_{ij}\Pi_{ij}D_{ij} : \Pi\mathbf 1=q,\ \Pi^{\!\top}\mathbf 1=\widehat p_t^v\}$ is \emph{linear} in the ground cost $D$, which is exactly what renders the dual \eqref{eq:w_dro_dual} a single-parameter (in $\lambda$) linear program. The ground cost $D_{ij}$ between two \emph{modes} is in turn the \emph{Wasserstein-2} (Bures) distance between their predicted Gaussian trajectory distributions,
\[
D_{ij}=\tfrac{1}{N_s}\sum_{k\le N_s}\!\Big(\|\mu_{i,k}^v-\mu_{j,k}^v\|^2+\mathcal B^2\!\big(\Sigma_{i,k}^v,\Sigma_{j,k}^v\big)\Big)^{1/2},
\]
where $\mathcal B$ is the Bures metric on covariance matrices. Hence the \emph{outer} ball is $W_1$ (chosen for the tractable linear dual and the concentration-calibrated radius), while the \emph{inner} mode-to-mode metric is $W_2$-Bures (the canonical closed-form distance between Gaussians). This resolves an apparent tension: a finite set of behavior modes is not \emph{a priori} a metric space, but embedding each mode through its dynamics into Bures--Wasserstein space supplies a genuine ground metric on which the outer $W_1$ ball is well defined. Using the true metric distance (rather than a squared cost) keeps $W_D$ a proper metric, so the radius $\rho$ retains its metric meaning.}
Let $n_{\mathrm{obs}}$ denote the number of mode observations available in the history window for obstacle class $C_v$.
We set
\[
\rho \;=\; \rho_0\!\left(1+\frac{\alpha}{\sqrt{n_{\mathrm{obs}}}}\right),
\]
where $\rho_0$ is a base radius (default $\rho_0=0.1$) and $\alpha$ is a user-set scaling coefficient (default $\alpha=1.0$).
This choice inflates the Wasserstein ball when observations are scarce and contracts as data accumulates, since $1/\sqrt{n_{\mathrm{obs}}}$ decays, yielding $\rho\to\rho_0$ in the well-observed regime.

\chg{The heuristic above adapts to $n_{\mathrm{obs}}$ but plateaus at $\rho_0$ and carries no coverage guarantee: nothing ensures the true mode law $p^\star$ lies inside the ball, which is the property the worst-case argument requires. We therefore calibrate $\rho$ as a genuine $W_D$ concentration radius. The empirical belief $\widehat p_t^v$ over $M$ modes, formed from $n$ observations, concentrates in total variation as $\Pr\!\bigl(\|\widehat p_t^v-p^\star\|_1\ge\epsilon\bigr)\le 2^M e^{-n\epsilon^2/2}$ (Devroye's simplex concentration inequality), so at target miscoverage $\beta$ the $L_1$ half-width is $\epsilon_n(\beta)=\sqrt{2\,(M\ln 2+\ln(1/\beta))/n}$. Since $W_D(q,p)\le\tfrac12\,\mathrm{diam}(D)\,\|q-p\|_1$ for any metric ground cost $D$, we set
\[
\rho_n(\beta)\;=\;\rho_{\min}+\tfrac12\,\mathrm{diam}(D)\,\epsilon_n(\beta),\qquad \mathrm{diam}(D):=\max_{i,j}D_{ij}.
\]
This radius is the principled choice for the reweighting on three counts. (i)~\emph{Metric consistency}: the ambiguity set is a $W_D$ ball, so its radius must be the matching $W_D$ concentration half-width; this is exactly what guarantees $p^\star\in\{q:W_D(q,\widehat p_t^v)\le\rho_n(\beta)\}$ with probability at least $1-\beta$, and hence $\langle q^\star,r\rangle\ge\langle p^\star,r\rangle$ at the same confidence---the reweighted expected violation upper-bounds the true one. (ii)~\emph{Statistical consistency}: $\rho_n(\beta)\to\rho_{\min}$ as $n\to\infty$, so the hedging vanishes and nominal scenario sampling is recovered once the belief is well estimated; the robustness costs conservatism only while the estimate is genuinely uncertain. (iii)~\emph{Geometry awareness}: it folds the transport diameter $\mathrm{diam}(D)$ of the ground cost in explicitly, rather than absorbing it into a base radius. We clamp $\rho_n(\beta)$ to $[\rho_{\min},\rho_{\max}]$, with $\rho_{\max}$ set below the mode-transport-collapse threshold so that $q^\star$ remains graded rather than degenerating to a single mode.}

We form $r$ using a conservative Gaussian clearance surrogate \cite{blackmore2012CCOPPW}. For mode $m$ and safe-horizon step $k$, let the predicted obstacle position satisfy $x_{t+k|t}^{v,m}\sim\mathcal N(\mu_{m,k}^v,\Sigma_{m,k}^v)$ and let $c_d(k)$ denote the center of ego disc $d$ at step $k$.
We then define a normalized straight line from the center of the ego disc to the obstacle $n_{m,k,d}:=(\mu_{m,k}^v-c_d(k))/\|\mu_{m,k}^v-c_d(k)\|$ and calculate the directional standard deviation for the obstacle $\sigma_{m,k,d}:=\sqrt{n_{m,k,d}^\top \Sigma_{m,k}^v\, n_{m,k,d}}$. 
Given a one-sided quantile $z_\alpha$, the per-step violation is
\[
r_{m,k,d} := \bigl[R + z_\alpha \sigma_{m,k,d} - \|\mu_{m,k}^v-c_d(k)\|\bigr]_+,
\]
where $R=r_d+r_{\mathrm{obs}}^v+r_{\mathrm{offset}}$ is the total safety radius and $[x]_+:=\max\{0,x\}$ is a functional taken over the quantity. We then set $r_m := \max_{k\le N_s}\max_{d\le D} r_{m,k,d}$.

\chg{The score $r_{m,k,d}$ is a \emph{surrogate}: it projects the two-dimensional collision violation onto the mean ego-to-obstacle direction $n_{m,k,d}$ and treats it as Gaussian with directional standard deviation $\sigma_{m,k,d}$. By Cauchy--Schwarz this projection upper-bounds the true Euclidean violation, so the surrogate is closed-form (no Monte Carlo) and conservative. Evaluating the quantile at the Bonferroni-corrected level $\alpha' = 1-(1-\alpha)/(N_sD)$ makes $\max_{k,d} r_{m,k,d}$ a valid upper bound on the \emph{joint}-horizon VaR of Euclidean collision: writing $V_{k,d}$ for the true violation, a union bound over the $N_sD$ per-step, per-disc events gives $\Pr[\max_{k,d} V_{k,d} > t] \le \sum_{k,d}\Pr[V_{k,d}>t] \le N_sD\,(1-\alpha') = 1-\alpha$. Two tighter but costlier variants are available. The \emph{proper} Bonferroni VaR keeps this union correction but replaces the linearised $\sigma$-surrogate with the true per-step Euclidean VaR, estimated by Monte Carlo on each step's marginal Gaussian; this removes the Cauchy--Schwarz slack while remaining a valid joint upper bound. The \emph{joint} VaR/CVaR is the exact risk measure of $V=\max_{k,d}[R-\|x_k-c_d(k)\|]_+$ over the correlated mode rollout, estimated by Monte Carlo over the whole horizon. We adopt the closed-form surrogate Bonferroni VaR by default for efficiency: it avoids any per-step Monte Carlo while retaining the joint-horizon guarantee.}

Note that $\widehat p_t^v$ is feasible for the Wasserstein ambiguity set
$\{q \in \Delta : W_D(q,\widehat p_t^v)\le \rho\}$ since
$W_D(\widehat p_t^v,\widehat p_t^v)=0\le \rho$.
Therefore, for the worst-case reweighting problem
$q^\star \in \arg\max_{W_D(q,\widehat p_t^v)\le \rho}\langle q, r\rangle$,
we have
$\langle q^\star, r\rangle \ge \langle \widehat p_t^v, r\rangle$.
That is, the WDRO produced $q^\star$ is at least as conservative as the nominal
distribution in terms of expected clearance violation. This means we can generate scenarios
according to $q^\star$ when forming the scenario constraints, because the reweighted dsitribtuion is within the orignal. 

Using the Kantorovich dual for the $W_1$-type objective, \eqref{eq:w_dro_modes} reduces to a 1D convex minimization in $\lambda\ge 0$:
\begin{equation}\label{eq:w_dro_dual}
\min_{\lambda\ge 0}\ \Bigl\{
\lambda\rho+\sum_{i=1}^M \widehat p_{t,i}^v\ \max_{j\in\mathcal M}\bigl(r_j-\lambda D_{ij}\bigr)
\Bigr\},
\end{equation}
which we solve via a 1D search \cite{gao2022distributionallyrobuststochasticoptimization}.
Here, $\lambda$ is the Lagrange multiplier (price per unit transport), and $\max_j(r_j-\lambda D_{ij})$ selects the best destination mode $j$ for mass starting at mode $i$ under the transport geometry.

We recover a candidate worst-case distribution by mapping each nominal mass at mode $i$ to
$j^\star(i)\in\operatorname*{arg\,max}_{j\in\mathcal M}(r_j-\lambda^\star D_{ij})$ and verify that the induced transport plan satisfies the radius constraint (up to numerical tolerance) before using the resulting weights for sampling.

\subsection{DRO scenario injection}

Beyond using WDRO to reweight the sampling distribution, we also use it to inject a small number of deterministic, mode-conditioned constraints into the SH-MPCC. Concretely, we select a set of $K$ modes and add the corresponding deterministic trajectory constraints directly to the optimization problem. We refer to the WDRO-based selection as WDRO-$K$. For comparison, we also consider two baseline selection rules that do not solve the WDRO problem: TopRisk-$K$ and DiverseRisk-$K$. Importantly, this injection step is not treated as additional i.i.d.\ scenario sampling. Instead, we interpret it as constraint tightening and add deterministic injected constraints which can only shrink the feasible set of the original scenario program. This therefore maintains the standard scenario-theoretic guarantee for the $S$ i.i.d.\ sampled constraints in like with the tehoretical requirements of \cite{Cal2006ScenarioApproach}.

\paragraph{WDRO-$K$ injection.}
WDRO-$K$ injects the $K$ modes receiving the largest mass under the WDRO worst-case distribution
$q^\star \in \Delta_M$, obtained by solving \eqref{eq:w_dro_modes} (via the dual \eqref{eq:w_dro_dual}).
Let $q^\star_m$ denote the weight assigned to mode $m\in\mathcal{M}$.
We select the index set
\[
\mathcal{I}_K \subset \mathcal{M}, \qquad |\mathcal{I}_K|=K, \qquad 
\mathcal{I}_K \in \operatorname*{arg\,max}_{|\mathcal{I}|=K}\ \sum_{m\in\mathcal{I}} q^\star_m,
\]
equivalently the $K$ largest components of $q^\star$, and inject the corresponding deterministic
representatives (e.g., mode-conditioned mean trajectories) as hard constraints. Unlike TopRisk-$K$,
WDRO-$K$ incorporates both the per-mode risk scores $r_m$ and the transport geometry induced by $D$
through the Wasserstein-ball constraint $W_D(q,\widehat p_t^v)\le \rho$, thereby emphasizing modes
that are dangerous and reachable under the Wasserstein cost structure.

\paragraph{TopRisk-$K$ injection.}
TopRisk-$K$ is a greedy risk-only heuristic. For each mode $m$, we compute a scalar risk score $r_m$ (e.g., the maximum clearance-surrogate violation over steps $k\le N_s$ and ego discs $d\le D$). We then select the $K$ modes with the largest risk scores and inject their deterministic representatives as trajectory constraints. This ensures the optimizer explicitly accounts for the most dangerous modes, but it may select redundant modes with nearly identical predicted trajectories.

\paragraph{DiverseRisk-$K$ injection.}
To reduce redundancy among injected modes, DiverseRisk-$K$ explicitly trades off risk against geometric diversity. Let $D(m,m')$ denote the pairwise transport-based distance between modes induced by the WDRO ground-cost matrix (e.g., $W_2$-Bures). The method initializes with the highest-risk mode and then greedily adds modes that are simultaneously high-risk and far from those already selected:
\begin{align}
m_1 &\in \arg\max_m r_m, \notag\\
m_{\ell} &\in \arg\max_{m\notin S_{\ell-1}}
\Bigl(r_m \cdot \min_{s\in S_{\ell-1}} D(m,s)\Bigr), \notag\\
S_{\ell} &= S_{\ell-1}\cup\{m_{\ell}\},
\qquad \ell=2,\dots,K. \notag
\end{align}
The resulting $K$ injected modes yield a set of hard constraints that better covers distinct failure patterns than TopRisk-$K$, while still prioritizing high-risk behaviors.

\paragraph{Constraint tightening.}
Let $\mathcal F_S$ denote the feasible set defined by the $S$ i.i.d.\ sampled scenario constraints, and let $\mathcal F_{S,\mathrm{inj}}$ denote the feasible set after adding injected deterministic constraints. Then $\mathcal F_{S,\mathrm{inj}}\subseteq \mathcal F_S$, so injection acts as a tightening heuristic and does not alter the scenario-theoretic guarantee on the sampled constraints.

%==============================================================================
% SECTION IV: METHODOLOGY
%==============================================================================

\section{Methodology}\label{methods}

We present a scenario-based MPCC framework for autonomous navigation under obstacle uncertainty with switching dynamics.
Collision avoidance is enforced through Safe Horizon-style scenario constraints \cite{degroot2023SBMotionBounded} using linearized half-space constraints.
To mitigate under-sampling of rare switching modes at practical scenario budgets, we introduce Wasserstein-DRO mode reweighting and a scenario-injection heuristic based on the worst-case distribution $q^\star$.

\subsection{Ego vehicle model}\label{sec:ego_model}

The ego vehicle uses a second-order unicycle model with contouring control. The ego state is
\begin{equation}\label{eq:ego_state}
x_t^{\mathrm{ego}}
=
\begin{bmatrix}
p_{x,t} & p_{y,t} & \theta_t & v_t & s_t
\end{bmatrix}^\top
\in \mathbb{R}^5,
\end{equation}
where $(p_{x,t}, p_{y,t})$ is position, $\theta_t$ is heading, $v_t$ is longitudinal speed, and $s_t$ is progress along a reference path $\Gamma$.
The control input is
\begin{equation}
u_t =
\begin{bmatrix}
a_t \\ \omega_t
\end{bmatrix}.
\end{equation}

The continuous-time dynamics are
\begin{subequations}\label{eq:ego_ct}
\begin{align}
\dot p_x &= v\cos\theta, &
\dot p_y &= v\sin\theta, \\
\dot \theta &= \omega, &
\dot v &= a, \\
\dot s &= v_s,
\end{align}
\end{subequations}
where $v_s$ denotes the progress rate along $\Gamma$ and $u_t$ is held constant over each integration interval of length $\Delta t$.
We discretize \eqref{eq:ego_ct} using one step of fourth-order Runge--Kutta (RK4):
\begin{subequations}\label{eq:ego_dyn_rk4}
\begin{align}
x_{t+1}^{\mathrm{ego}} &= x_t^{\mathrm{ego}} + \frac{\Delta t}{6}\left(k_1 + 2k_2 + 2k_3 + k_4\right),\\
k_1 &= f\!\left(x_t^{\mathrm{ego}},u_t\right),\quad
k_2 = f\!\left(x_t^{\mathrm{ego}} + \frac{\Delta t}{2}k_1,\ u_t\right),\\
k_3 &= f\!\left(x_t^{\mathrm{ego}} + \frac{\Delta t}{2}k_2,\ u_t\right), \\
k_4 & = f\!\left(x_t^{\mathrm{ego}} + \Delta t\,k_3,\ u_t\right),
\end{align}
\end{subequations}
with
\[
f(x,u)=
\begin{bmatrix}
v\cos\theta\\
v\sin\theta\\
\omega\\
a\\
v_s
\end{bmatrix}.
\]

\subsection{Model Predictive Contouring Control cost}

Given the reference path $\Gamma$, the MPCC cost penalizes two geometric errors at each horizon step $k$:
\begin{itemize}
    \item \textbf{Contouring error} $e_c(k)$: the signed distance from the predicted ego position to the closest point on $\Gamma$ at progress $s_{t+k|t}$.
    \item \textbf{Lag error} $e_l(k)$: the along-path deviation between $s_{t+k|t}$ and the desired progress.
\end{itemize}
The MPCC stage cost at step $k$ is
\begin{align}
    \ell_{\mathrm{MPCC}}(k) = w_c\, e_c(k)^2 + w_l\, e_l(k)^2 + \notag \\ w_a\, a_{t+k|t}^2 + w_\omega\, \omega_{t+k|t}^2 + w_v (v_{t+k|t} - v_{\mathrm{ref}})^2, \label{eq:mpcc_cost}
\end{align}
where $w_c, w_l, w_a, w_\omega, w_v > 0$ are user-specified weights and $v_{\mathrm{ref}}$ is a reference velocity. The reference path $\Gamma$ is represented as a cubic spline, and the contouring/lag errors are computed from the spline evaluation at $s_{t+k|t}$.

\subsection{Multi-disc robot model}\label{sec:multidisc}

The robot footprint is approximated by $D$ discs of radius $r_d$ centered at offsets $\{c_1,\dots,c_D\}$ along the robot's longitudinal axis. The center of disc $d$ at horizon step $k$ is
\begin{equation}\label{eq:disc_center}
p_d(k) = \begin{bmatrix} p_x(k) + c_d \cos\theta(k) \\ p_y(k) + c_d \sin\theta(k) \end{bmatrix},
\end{equation}
where we abbreviate $p_x(k) := p_{x,t+k|t}$, etc. Multi-disc models provide tighter footprint coverage than a single bounding circle, which is critical for non-circular robots.

\subsection{Obstacle model, uncertainty, and switching}

Obstacle agents are indexed by $v \in \{1,\dots,V\}$. For each obstacle $v$, a motion predictor provides a predicted distribution over future positions. We assume access to a fast base predictor from which particles can be sampled, e.g., Gaussian or Gaussian mixture models.

We assume each obstacle $v$ belongs to a class $\mathcal{C}(v)$ with a finite mode set $\mathcal{M}_{\mathcal{C}} = \{1,\dots,M_{\mathcal{C}}\}$, where mode $m$ specifies a stochastic motion model
\begin{equation}\label{eq:obs_dyn}
x_{t+1}^v = f_m(x_t^v, \omega_t^v), \qquad \omega_t^v \sim \mathbb{P}_{\omega,m}.
\end{equation}
The obstacles follow the identical state formulation to the ego vehicle as descibed in Section~\ref{sec:cc_prelim}. Here, $\mathbb{P}_{\omega,m}$ denotes the distribution of all uncertainties associted with given mode $m$.
The mode evolution process is unknown and may be non-Markovian.

\subsection{Linearized half-space collision constraints}\label{sec:linearized}

Collision avoidance between each robot disc $d$ and obstacle $v$ at horizon step $k$
is enforced using a linearized half-space constraint for each scenario $i$.
Let $x^{v,(i)}_{t+k|t}\in\mathbb{R}^2$ denote the sampled obstacle position at step $k$ in
scenario $i$. Define the outward normal from disc $d$ to obstacle $v$:
\begin{equation}\label{eq:normal}
\hat{n}_{d,v,k}^{(i)} =
\frac{x^{v,(i)}_{t+k|t} - p_d(k)}
{\max\!\left(\|x^{v,(i)}_{t+k|t} - p_d(k)\|, \varepsilon_{\min}\right)},
\end{equation}
where $\varepsilon_{\min} > 0$ guards against division by zero. The half-space
constraint is
\begin{equation}\label{eq:halfspace}
\hat{n}_{d,v,k}^{(i)\top} p_d(k)
\le
\hat{n}_{d,v,k}^{(i)\top} x^{v,(i)}_{t+k|t}
- (r_d + r_{\mathrm{obs}}^v + r_{\mathrm{offset}}),
\end{equation}
where $r_{\mathrm{offset}}\ge 0$ is an additional safety margin.

\chg{Before solving, we reduce the $S\!\times\!N_s\!\times\!D$ half-spaces by \emph{geometric dominance pruning} (de Groot et al.\ \cite{degroot2023SBMotionBounded}, Definition 2). A scenario $\delta^{(i)}$ is redundant when its feasible set contains that of another scenario, $\Theta_{\delta^{(j)}}\subseteq\Theta_{\delta^{(i)}}$: satisfying the more-constraining $\delta^{(j)}$ then automatically satisfies $\delta^{(i)}$, so $\delta^{(i)}$ can be removed without changing the optimizer's solution. Rather than approximate this by a distance-plus-alignment heuristic, we certify each pairwise implication \emph{exactly}. For a matched step $k$, ego disc $d$, and obstacle, the two linearized constraints are half-spaces $H_q=\{p:a_q^\top p\ge b_q\}$ with $a_q$ the unit obstacle-to-ego normal and $b_q=a_q^\top x_{\text{obs}}^{(q)}+R$. We restrict attention to a bounded \emph{reachable} ball $\mathcal B=\{p:\lVert p-\bar p_{k,d}\rVert\le\rho\}$ about the reference disc center and test whether $H_j\cap\mathcal B\subseteq H_i$ by solving, in closed form, $\min_{p\in\mathcal B,\;a_j^\top p\ge b_j}\,(a_i^\top p-b_i)\ge 0$ (a linear program over a ball intersected with a half-space). Scenario $\delta^{(j)}$ dominates $\delta^{(i)}$ only when this implication holds for \emph{every} step, disc, and obstacle, when the two scenarios constrain an identical obstacle set over the full horizon, and $\rho$ over-estimates the true reachable displacement. This is a \emph{conservative sufficient condition} for $\Theta_{\delta^{(j)}}\subseteq\Theta_{\delta^{(i)}}$: it may leave some genuine redundancies un-pruned, but it never certifies a containment that fails, so an active (non-redundant) constraint is never discarded and the reduction is non-distorting by construction. Because the certified relation is true containment on a fixed ball it is transitive, so the greedy pairwise elimination is order-independent for safety; DRO-injected worst-case scenarios are additionally never pruned. In the limit $\rho\to\infty$ the test reduces to exact half-space nesting (identical normals, closer obstacle), so with a large $\rho$ the reduction removes only near-collinear duplicate scenarios---exactly the redundancies for which non-distortion is guaranteed.}

\subsection{Simulation Environments}
\label{sec:envs}

We evaluate the proposed SH-MPCC with WDRO reweighting and mode-injection in two canonical
interaction scenarios designed to stress different failure modes of probabilistic prediction and planning.
In both environments, the ego vehicle initializes at $(0,0)$ and tracks a fixed $25$\,m reference path given
by an S-curve
\[
y(x) \;=\; 3\sin\!\Bigl(\tfrac{2\pi x}{L}\Bigr), \qquad L=25\ \mathrm{m},
\]
and we parameterize locations along the path by arc length to specify obstacle placement consistently.

\paragraph{Oncoming environment.}
The obstacle is initialized directly on the ego path at $60\%$ of the path arc length (approximately
$15$\,m ahead of the ego) and moves head-on toward the ego at roughly $1$\,m/s, yielding an
adversarial, near-worst-case interaction on a direct collision course. The obstacle behavior is modeled
as a multimodal predictor with a small discrete set of maneuver modes: four common modes
(\emph{constant velocity}, \emph{turn left}, \emph{turn right}, \emph{decelerate}) together with a rare
\emph{lane-change} mode. Under the nominal predictor (no WDRO reweighting and no injection),
this environment exhibits a high base collision rate of approximately $67\%$, making it the most
challenging setting and a stringent test of conservativeness and safe-horizon constraint tightening.

\paragraph{Intersection environment.}
To probe performance under crossing interactions and mode ambiguity, we consider an intersection-like
setup on the same S-curve reference path. Here the obstacle is initialized $2$\,m laterally off the path at
$40\%$ of the path arc length and moves approximately perpendicular to the ego path at $\sim 1$\,m/s,
mimicking a vehicle crossing the ego lane (a T-intersection geometry). In this setting, the ego must
anticipate whether the obstacle will \emph{cross}, \emph{stop}, or \emph{turn} as it enters the interaction
region. The nominal base collision rate is comparatively low ($\approx 4$--$7\%$), but performance is more
sensitive to mode coverage: the crossing angle and short time-to-interaction make prediction errors
and missed modes more consequential, providing a complementary stress test to the oncoming case.

\subsection{Scenario generation}\label{sec:scenario_gen}

A scenario consists of obstacle trajectories for all obstacles over the safe horizon.
For each obstacle $v$, let the base predictor define a distribution over the stacked
vector $\xi_t^v \in \mathbb{R}^{2N_s}$ (concatenated obstacle positions across the safe horizon).
For each rollout, we generate scenarios by first sampling a mode $m$ and then sampling
trajectory noise from the base predictor conditioned on that mode.
In base SH-MPCC, modes are sampled from $\widehat{p}_t^v$. In WDRO-sampling tests, modes are sampled from $q^\star$. In injection variants, the $S$ sampled scenarios are drawn from the chosen sampling distribution and $K$ additional deterministic mode constraints are injected.

\subsection{Safe Horizon scenario optimization}\label{sec:safe_horizon}

We enforce collision avoidance constraints only over a safe horizon $N_s \le N$
\cite{degroot2023SBMotionBounded}. All collision constraints and mode-risk computations
are truncated to $N_s$, while the control horizon for optimization remains $N$.
This keeps the scenario program manageable (we certify safety only on the truncated
horizon) while still planning controls over the full horizon. All risk scores $r$
being used in equations \eqref{eq:w_dro_modes} and \eqref{eq:w_dro_dual} are computed using only timesteps
$k\le N_s$, aligning DRO reweighting with the certified risk horizon.

\chg{The safe horizon $N_s$ is not a free parameter: it is bounded by the scenario budget through the sample complexity of the scenario approach \cite{Cal2006ScenarioApproach,Ca2018ScenarioTheory}. Guaranteeing the chance constraint at violation level $\epsilon$ with confidence $1-\beta$ requires $S\gtrsim \tfrac{2}{\epsilon}\ln\tfrac1\beta+2\bar n+\tfrac{2\bar n}{\epsilon}\ln\tfrac2\epsilon$ scenarios, where $\bar n$ is the support rank (the number of scenario constraints simultaneously active at the optimum). Because support-set selection caps $\bar n$ at a fixed value independent of the horizon length, this budget is \emph{horizon-independent}: a given $S$ either certifies the entire horizon or none of it. At representative levels ($\epsilon=0.05$, $\beta=0.01$, $\bar n=5$) the certified budget is $S\approx 10^3$ scenarios---an order of magnitude beyond real-time budgets---so the certified safe horizon collapses to a short prefix $N_s$; scaling $N_s$ with the available budget instead trades the formal certificate for a longer effective planning horizon. This truncation sharpens the motivation for WDRO reweighting: a short $N_s$ evaluated with few certified scenarios is precisely the regime in which an unhedged nominal sampler is most likely to omit a rare high-risk mode, which is what $q^\star$ counteracts \cite{degroot2023SBMotionBounded}.}

With the assumption of no disturbances, the optimization follows:

\begin{subequations}\label{eq:scenario_opt}
\begin{align}
\min_{U_t \in \mathbb{U}^N}\quad &
\sum_{k=0}^{N-1} \ell_{\mathrm{MPCC}}(k) \label{eq:obj}\\
\text{s.t.}\quad &
x_{t+k+1|t}^{\mathrm{ego}} = f^{\mathrm{ego}}(x_{t+k|t}^{\mathrm{ego}}, u_{t+k|t}),
\quad x_{t|t}^{\mathrm{ego}} = x_t^{\mathrm{ego}}, \label{eq:dyn_constr}\\
& \hat{n}_{d,v,k}^{(i)\top} p_d(k)
\le
\hat{n}_{d,v,k}^{(i)\top} x^{v,(i)}_{t+k|t}
- (r_d + r_{\mathrm{obs}}^v + r_{\mathrm{offset}}), \label{eq:coll_constr}\\
& \quad \forall\, d=1,\dots,D,\ v=1,\dots,V \notag \\
& \quad \forall\, k=1,\dots,N_s,\ i=1,\dots,S_{\mathrm{tot}},
\nonumber
\end{align}
\end{subequations}

\subsection{Algorithm}\label{sec:algorithm}

Algorithm~\ref{alg:full_method} summarizes the full method, including WDRO reweighting.
\begin{algorithm}[ht!]
\caption{Scenario Safe Horizon MPCC with WDRO reweighting}\label{alg:full_method}
\begin{algorithmic}[1]
\Require Ego state $x_t^{\mathrm{ego}}$, obstacles $\{x_t^v\}_{v=1}^V$, horizon $N$, safe horizon $N_s$, risk $(\epsilon,\beta)$, discs $D$
\State Compute scenario sample size $S$ from scenario theory (implementation bound)
\For{$t=0,1,2,\dots$}
    \For{each obstacle $v=1,\dots,V$}
        \State Obtain nominal mode belief $\widehat{p}_t^v$ and mode-conditioned predictions
       \State Compute $q^\star$ by Wasserstein-DRO reweighting in \eqref{eq:w_dro_modes}
(equivalently dual \eqref{eq:w_dro_dual}) using ground cost $D$
    \EndFor
    \State Sample $S$ i.i.d.\ scenarios using modes $q^\star$ and base predictor noise
    \State Build half-space constraints over $k=1,\dots,N_s$ via \eqref{eq:normal}--\eqref{eq:halfspace}
    \State Solve scenario MPCC \eqref{eq:scenario_opt} to obtain $U_t^\star$
    \State Apply $u_{t|t}^\star$
\EndFor
\end{algorithmic}
\end{algorithm}

%==============================================================================
% SECTION V: RESULTS
%==============================================================================

\section{Results}\label{results}

We evaluate safety (collision rate), mode coverage (missed-mode rate), clearance, and runtime in two environments (Oncoming and Intersection). We focus on four empirical questions: (Q1) does WDRO-sampling improve mode coverage at fixed scenario budget? (Q2) does WDRO injection improve safety, and when does WDRO out perform simpler TopRisk and DiverseRisk injection? (Q3) is performance robust against short mode history lengths? (Q4) how sensitive are results to the transport ground cost choice?

\paragraph{Methods compared.}
We report base SH-MPCC (sample modes from $\widehat p_t^v$),
WDRO-sampling (sample modes from $q^\star$),
TopRisk and Softmax-risk baselines,
and three injection policies for injection budget $K$: WDRO-$K$, TopRisk-$K$, and DiverseRisk-$K$.

%------------------------------------------------------------------------------
\subsection{Mode-set injection under finite budgets (D1)}\label{sec:res_d1}
We first study \emph{mode-set injection}: injecting $K\ge 1$ deterministic scenarios corresponding to a selected set of modes. Figure~\ref{fig:figA_selection} compares WDRO-$K$, TopRisk-$K$, and DiverseRisk-$K$. In the Oncoming situation, WDRO-$K$ achieves the lowest collision rates and provides its largest gains when $K>1$, indicating that transport structure helps avoid redundant injections. In Intersection, overall collision rates are lower and confidence intervals overlap more; here the diversity baseline is competitive at larger $K$. These results suggest WDRO is most beneficial when a finite injection budget must represent multiple distinct risky behaviors rather than only the single most dangerous mode. 

\paragraph{Effect size and interpretability.}
DiverseRisk-$K$ already leverages the same transport ground cost \(D\) by formulation, and empirically its performance is close to WDRO-$K$ (differences are typically \(<1\) percentage point and sometimes change sign across \(K\)).
With finite rollouts, these small effects are difficult to separate statistically, so we avoid claiming a consistent empirical advantage of WDRO-$K$ over DiverseRisk-$K$.
Instead, our results support the broader point that transport geometry is the key ingredient for improving injection beyond risk-only baselines, while WDRO offers a principled selection rule derived from a worst-case expectation objective and a theoretical backing.

\chg{\begin{figure*}[t]
  \centering
  \includegraphics[width=\linewidth]{paper_figures/fig_a_selection.png}
  \caption{D1: WDRO-$K$ vs TopRisk-$K$ vs DiverseRisk-$K$ ($K=1,2,3,4$) in Oncoming and Intersection. WDRO provides its largest gains when $K>1$, consistent with improved mode-set selection under finite budgets.}
  \label{fig:figA_selection}
\end{figure*}}

%------------------------------------------------------------------------------
\subsection{WDRO-sampling vs risk-softmax sampling (D2)}\label{sec:res_d2}
We next evaluate WDRO-sampling against risk-softmax sampling $p_\tau(m)\propto \exp(\tau r_m)$ for multiple configurations. Figure~\ref{fig:figB_sampling} shows that WDRO-sampling reduces missed-mode rates by approximately 8--10 percentage points relative to all softmax variants in both environments, while achieving comparable collision rates. This indicates WDRO redistributes probability mass in a structured way that improves finite-sample coverage under switching, rather than merely amplifying the highest-risk mode.

\chg{\begin{figure*}[t]
  \centering
  \includegraphics[width=\linewidth]{paper_figures/fig_b_sampling.png}
  \caption{D2: WDRO-sampling vs risk-softmax sampling (multiple temperatures) in Oncoming and Intersection. WDRO-sampling yields substantially better mode coverage (lower missed-mode rate) at comparable collision rates.}
  \label{fig:figB_sampling}
\end{figure*}}

%------------------------------------------------------------------------------
\subsection{Mode-history length comparison (D3)}\label{sec:res_d3}
We vary the maximum history length used to estimate the nominal belief $\widehat p_t^v$ and to set the adaptive radius $\rho$. Figure~\ref{fig:figC_observation} shows WDRO-injection is the most stable method across history lengths in Oncoming, while TopRisk injection degrades as history grows. In Intersection, collision rates are uniformly low, but the mode-coverage advantage persists: injection policies achieve lower missed-mode rates than base SH-MPCC. These results support uncertainty-adaptive ambiguity sets for switching environments.

\chg{\begin{figure*}[t]
  \centering
  \includegraphics[width=\linewidth]{paper_figures/fig_c_observation.png}
  \caption{D3: History-length sweep in Oncoming and Intersection. WDRO-injection is most stable in the hardest environment, and coverage benefits persist across regimes.}
  \label{fig:figC_observation}
\end{figure*}}

%------------------------------------------------------------------------------
\subsection{Ground-cost sensitivity (D4)}\label{sec:res_d4}
We evaluate sensitivity to the transport ground cost $D$ using W2--Bures, Euclidean-mean, and zero--one costs. Figure~\ref{fig:figD_ground_cost} and Table~\ref{tab:d4_ground_cost} show performance is broadly consistent across costs for WDRO injection, while WDRO-sampling is mildly degraded under the zero--one cost. Overall, WDRO does not hinge on a fragile cost choice, which is favorable for practitioners.

\chg{\begin{figure}[t]
  \centering
  \includegraphics[width=\linewidth]{paper_figures/fig_d_ground_cost.png}
  \caption{D4: Ground-cost ablation. Injection performance is similar across costs; WDRO-sampling is mildly worse under the zero--one cost.}
  \label{fig:figD_ground_cost}
\end{figure}}

\begin{table}[t]
\centering
\caption{Collision rate (Coll.\%) and Missed Mode rate (MM\%) across ground costs.}
\label{tab:d4_ground_cost}
\resizebox{\columnwidth}{!}{%
\begin{tabular}{lcccccc}
\toprule
Method & \multicolumn{2}{c}{W2-Bures} & \multicolumn{2}{c}{Zero-One} & \multicolumn{2}{c}{Euclidean-Mean} \\
 & Coll.\% & MM\% & Coll.\% & MM\% & Coll.\% & MM\% \\
\midrule
WDRO-K2 & 13.7 & 23.8 & 14.1 & 23.4 & 14.5 & 23.3 \\
WDRO-K3 & 15.3 & 24.1 & 15.3 & 23.3 & 13.9 & 23.1 \\
WDRO-sampling & 15.9 & 27.7 & 17.7 & 27.5 & 15.8 & 27.1 \\
\bottomrule
\end{tabular}
%
}
\end{table}

\subsection{Takeaway}
Overall, the results indicate that WDRO yields several practical benefits. 
First, WDRO-based reweighting improves mode coverage---reflected in lower missed-mode rates---relative to baselines.
Second, WDRO is particularly helpful when used for mode selection in constraint injection: WDRO-$K$ injection is most effective when multiple modes are injected ($K>1$), where it can capture distinct high-risk behaviors that would otherwise be missed.
Third, the uncertainty-adaptive Wasserstein radius stabilizes performance under limited mode observations, and the qualitative ranking of methods is largely consistent across ground-cost choices, suggesting limited sensitivity to the specific cost metric and supporting practical deployment.

%==============================================================================
% SECTION VI: CONCLUSION
%==============================================================================

\section{Conclusion}\label{conclusion}

We studied Safe Horizon MPCC under switching obstacle dynamics and addressed a key practical failure mode of sampling based techniques. Namely, finite sample budgets can under-represent rare but safety-critical behaviors. We introduced a Wasserstein-DRO mode reweighting module over discrete motion modes and evaluated two complementary uses: WDRO-sampling and mode-based scenario injection.
WDRO-sampling consistently improved mode coverage under switching, reducing missed-mode rates by approximately 8--10 percentage points relative to baselines at comparable collision rates in two realistic environments.
For deterministic injection, greedy TopRisk injection is often competitive when only a single scenario is injected ($K=1$), but WDRO provides its largest gains when finite budgets must represent multiple distinct risky behaviors ($K>1$), improving injected mode-set selection relative to baselines.
A comparison across observation counts showed WDRO-injection is the most stable method as mode beliefs become uncertain, and ground-cost ablations indicated performance is not sensitive to the exact transport cost choice.
Overall, Wasserstein-DRO mode reweighting and lightweight injection provide a practical, low-overhead way to improve finite-sample coverage and safety in SH-MPCC under switching dynamics. Furture work will adapt these techniques to obstacles with non-discrete modes where DRO's geometric properties may further shine.

\bibliographystyle{IEEEtran}
\bibliography{tex}
\end{document}
```
